\documentclass[twoside]{article}
\usepackage[
  a5paper,
  margin=0.5in,
  includehead,
  headsep=4pt,
  headheight=13pt
]{geometry}
\usepackage{gregosheet}
\usepackage{gregosheet-psalm}
\usepackage{lyluatex}
\usepackage{xcolor}
\input{styles}

\begin{document}

\parttitle[ÉE 686]{Introitus}
\begin{gregosheet}
  \addinline{<X-3r-3---3--:-34t-tT4rR3r-rR3-:-34t-tT4-5-5uU6h445zZ5T4-3--34tT4R3r-rR3-,-tT4-tT4R3-5---4t-rR3-3-:-tT4-5--5uU6h445zZ5T4-34tT4R3r-rR3-,,-¨-----3r-rR3-4t--[---------------4--5-,-3--4--[--------------5---5--4--6--4--5--,-3--4t-[------------------------------3--4t-4--3--3-.}
            {   Re-qui-em * ae--ter-----nam   do--na  e-i(s),         Do-mi-------ne.   Et  lux   per-pe-tu--a   lu--ce-at            e--------i(s).  <Ps.> Te de--cet hymnus,_Deus_in Si-on, et ti-bi_reddetur_vo-tum in Ie-ru-sa-lem, ex-au-di_orationem_meam,_ad_te_omnis ca-ro ve-ni-et.}
\end{gregosheet}

\parttitle[GR 94*]{Kyrie}
\begin{gregosheet}
  \addinline{<X-3r5z-5--tT4x-:-tT4R3E2e-4-3-3y-,,-3r5z--tT4x-:-tT4R3E2e-4--3-3y-?-3r5z-5--tT4x-:-tT4R3E2e-4--3-3y-,,-7--3--3yuU67iuU6Z5T4x-:-tT4R3E2e-4--3-3y-.}
            {   Ky---ri-e    * e-le-i-son.        Chris-te   * e--------le-i-son.  Ky---ri-e    * e--------le-i-son. Ky-ri-e               * e--------le-i-son.}
\end{gregosheet}

\parttitle[ÉE 688]{Graduále}
\begin{gregosheet}
  \addinline{<-0---0-0--0---0---1---0---0e-3--:-3---3--3---1-2e----4--,-3----3--1---0----1e-3-.}
            {  Bár a ha-lál völ-gyé-ben já-rok, nem fé-lek a rossz-tól, mert te ott vagy vé-lem.}
\end{gregosheet}

\parttitle[GH 565]{Alleluja}
\begin{gregosheet}
  \addinline{<-34t-tT456uU6-5u-7iI7j5T4-,,-77-4tT4R3-=-56uuJ4uj5z5tT4-,,-¦-6-7-5-4-,,-77-4tT4R3-56uuJ4uj5z5tT4-.}
            {  Al--le-------lu-ja.      *}
\end{gregosheet}

\parttitle[ÉE 342]{Felajánlási ének}
\begin{gregosheet}
  \addinline{<ô-2--2--2---2---1-----2--3--4x-?-4-----4---5----6---5--4---3--2y-,-¨-----6----6---5----5---6--5--4--3y-,-3--4--5---6----5--4-3--2y-.}
            {   Se-re-gek-nek szent Is-te-ne,  menny-nek föld-nek Te-rem-tő-je!   <R)> Jöjj el, jöjj el, én Jé-zu-som, üd-vö-zít-sél vég-ó-rá-mon.}
\end{gregosheet}
\begin{hymnstanzas}
  2. Bűneimnek sokaságát\\
  Uram, ne nézd gonoszságát! R.
\stanzabreak
  3. Tekintsd lelki sebeimet,\\
  hozd gyógyító kegyelmedet! R.
\stanzabreak
  4. Néked lelkem fölajánlom,\\
  testemet a földnek szánom. R.
\stanzabreak
  5. Föltámad majd onnan egykor,\\
  az utolsó ítéletkor. R.
\end{hymnstanzas}

\parttitle[ÉE 448]{Sanctus}
\begin{gregosheet}
  \addinline{<X-5-----tT4---:-5-----tT4-:-3-----4-----[--------------5--4--5-,-3--4---[-------------4-6----5----5--4-tT4-3--,-3--4----5--5-5--5---4---5--,-34t-5---/-5-5--5---4--6--5--4tT4-3--,-1--3r---4-!-4-4--5---rR3-3r-.}
            {   Szent vagy, * szent vagy, szent vagy, mindenség_ura, Is-te-ne! Di-cső-séged_betölti a meny-nyet és a föl-det, Ho-zsan-na a ma-gas-ság-ban! Ál--dott, a-ki jön az Úr ne-vé---ben! Ho-zsan-na  a ma-gas-ság-ban!}
\end{gregosheet}

\parttitle[ÉE 449]{Agnus Dei}
\begin{gregosheet}
  \addinline{<-5--5---5--rR3-4---:--4--[-------------3---4--eE2-1-,-3--4---5----4---4---,,-5--5---5--rR3-4---:-4--[-------------3---4--eE2-1-,-3--4---5----4---4---,,-5--5---5--rR3-4---:--4--[-------------3---4--eE2-1-,-3---4--5----rR3E2-2-.}
            {  Is-ten Bá-rá--nya, * te elveszed_a_vi-lág bű-ne--it, ir-gal-mazz ne--künk!  Is-ten Bá-rá-nya, * te elveszed_a_vi-lág bű-ne--it, ir-gal-mazz ne--künk!  Is-ten Bá-rá-nya, *  te elveszed_a_vi-lág bű-ne-it,  adj ne-künk bé----két!}
\end{gregosheet}

\parttitle[ÉE 162]{Áldozási ének}
\begin{hymnstanzas}
  1. Üdvözlégy, Krisztusnak\\
  drága, szent teste!\\
  Áldott légy, Jézusnak\\
  piros, szent vére!
\stanzabreak
  R) Ó, mily áldott vagy, oltáriszentség,\\
  angyali drága vendégség!
\stanzabreak
  2. Üdvözlégy, nagy szentség,\\
  élet adója!\\
  Mi bűnös lelkünknek\\
  megorvoslója! R.
\stanzabreak
  3. Üdvözlégy, nagy szentség,\\
  világ váltsága!\\
  Emberi nemzetnek\\
  fő boldogsága! R.
\end{hymnstanzas}
\centerstanza{%
4. Üdvözlégy, nagy szentség,\\
lelkünk vendége!\\
Mennyei karoknak\\
gyönyörűsége! R.
}

\parttitle[ÉE 690]{Kommunió}
\begin{gregosheet}
  \addinline{<-5---rR3-4t--4-:-5--7--6--7-5---4--3r-4---,-5---7---6----7--5--6--7--5---4tT4-:-3---3-4t-5--4-,,-¨----4--5---7--7--7---7---7--7--6-7---8--7--7---,-5--7---7---7--7--7-7--6--7--5-4-?,-5---7---6----7--5-.}
            {  Lux ae--ter-na  lu-ce-at e-is, Do-mi-ne, * cum san-ctis tu-is in ae-ter-num,   qui-a pi-us es!  <V.> Re-qui-em ae-ter-nam do-na e-is, Do-mi-ne, * et lux per-pe-tu-a lu-ce-at e-is!  Cum san-ctis tu-is.}
\end{gregosheet}

%\begin{gregosheet}
%  \addinline[tone={1. lá tónus}]{<X-1t-5u-5----:-[------------------tT45uj5T4-4--:-5-4--3---3---3---:-2---3--4---3---2--1--0q-1--,-1--3----2---1-1--1---1---2----3----4tT4-3--:-3---2---3---4----eE2-1---1-,,-[-4-3-rR3-1-.}
%                                {   Én va-gyok * a_feltámadás_és_az é---------let, a-ki ben-nem hisz, még ha meg-hal is él-ni fog; és mind-az, a-ki úgy él, hogy hisz ben--nem, nem hal meg mind-ö---rök-ké.}
%\end{gregosheet}
%\begin{psalmsheet}[tone=1]
%  \psalmtext{
%    1. A mélységekből kiáltok hozzád, ó Uram, * Uram, hallgasd meg az én szómat!\\
%    2. Legyen a te füled figyelmes * az én könyörgésemnek szavára! ANT. ---
%
%    3. Ha a gonoszságokat számon tartod, Uram, * Uram, ki áll meg akkor teelőtted?\\
%    4. De nálad kegyelem vagyon, * és a te törvényed miatt reménykedem benned, Uram. ANT.
%}
%\end{psalmsheet}

\parttitle[ÉE 338]{Kivonulási ének}
\begin{gregosheet}
  \addinline{<XB-4--4----5--6--7--6--5--4x--,-7--6----5--4---3--5--4---3y--,-6-5--4--3--4-3--2---1í-,-4-----4---5-4---3---2---1--0í-.}
            {    Le-gyen ez az ál-do-za-tunk, me-lyet né-ked be-mu-tat-tunk, e-rő-sí-tő a ha-lál-ra   s_lel-kek ö-rök meg-vál-tá-sa.}
\end{gregosheet}

\centerstanza{%
Kik előttünk már elmentek,\\
lássák meg a fényességet,\\
s erejéből e szentségnek\\
nyerjék el az üdvösséget!
}

\end{document}